\subsection{问题一的建模与求解：一次往返场展开}

本问先回答条纹直接给出了什么量，再建立它与外延层厚度的联系。真实碳化硅光谱的周期等效量为$19.54\ \mu\mathrm{m}$，换算几何厚度还需光学参数。“一次往返场展开”是指：把外延层上表面的直接反射场，与光进入外延层、到达衬底界面后返回的首束场相加；光再次往返形成的更高阶返回不放入本问主模型。
% src:求解/问题1/结果/模型验证.json:真实附件诊断.表观光学厚度乘积_dq_微米

两谱有效周期约为$254.85$和$256.99\ \mathrm{cm}^{-1}$，以上述周期的中位数按第\ref{sec:thickness-quantities}节的定义折算。该实测量保留有限窗口内的色散和界面相位贡献，局部近似的条件由式\eqref{eq:period-equivalent}给出。
% src:求解/问题1/结果/模型验证.json:真实附件诊断.窗口_每厘米;真实附件诊断.周期投影诊断.0.周期_每厘米;真实附件诊断.周期投影诊断.1.周期_每厘米;真实附件诊断.表观光学厚度乘积_dq_微米

先限定为两束截断场，才能把角度、色散和偏振对条纹的作用逐项放入同一表达式，并把后续往返的收益单独检验；若直接引入高阶项，表面反射与首个返回的作用会失去清晰边界。两束的来源和真实数据的光学参数缺口将在分析与准备中交代，随后由界面导纳写出色散相位及其峰距关系。
